% =====================================================================
%  preface.tex —— 前言（原文 p.8–9）
% =====================================================================
\chapter*{前言}
\addcontentsline{toc}{chapter}{前言}
\markboth{前言}{前言}

% [原文 p.8]
本讲义源自作者近年为清华大学开设的“量子力学”与“量子力学的数学方法”两门课程。它意在同时给出量子力学的基本物理与现代量子化数学理论的引论。

第一部分偏重物理，旨在阐明量子力学的基础。第 1 章讲解传统的希尔伯特空间与算子进路下的量子力学基本原理。这一形式体系以酉变换与测量为语言，为描述量子系统的演化建立了框架。第 2 章讲解量子力学的路径积分进路，它以概率式的路径来刻画粒子的量子动力学。这一形式体系为计算跃迁振幅、理解量子现象提供了强有力的工具，并在现代量子场论中被广泛推广与发展。这两种等价的形式体系彼此深度交织，展示了量子力学这一激动人心的发展的不同侧面。

% [原文 p.9]
第二部分转向数学的观点，聚焦于量子算子的结构性基础。第 3 章讲述经由伪微分算子的量子化——一种把相空间上以函数表示的经典可观测量，系统地转换为作用在希尔伯特空间上的量子算子的方法。这一进路为沟通经典力学与量子力学提供了严格的框架。第 4 章介绍希尔伯特空间上算子的谱理论，它用谱数据来刻画量子可观测量。无论是诠释测量结果、对系统状态分类，还是分析动力学性质，都要用到它。这两个相互关联的理论，体现了数学分析如何把量子力学的概念基础形式化并加以扩展。

第三部分继续采取数学的观点，强调量子化的现代几何与代数进路。第 5 章讨论量子力学的相空间表述与几何量子化理论。第 6 章聚焦于形变量子化，这是一个代数框架，其中量子可观测量是经典对应物的形式形变。在这一表述中，普朗克常数 $\hbar$ 充当一个形式参数，它编码了一族非交换代数。

本讲义所呈现的材料大体上是标准的。除 §2.8.5 与 §6.4.3 两节（讨论拓扑量子力学与形变量子化中迹映射之间的联系）之外，作者不主张任何原创性。本工作取资于若干经典文献，见参考文献。

我非常感谢求真书院的“笔记记录员计划”（Notes Taker Program），它促成了本讲义的产生并给予支持。本讲义的初稿由杨鹏录入，我对他的出色工作深表谢意。

\vspace{1.2em}
\begin{flushright}
静斋\\
2025 年 12 月
\end{flushright}

\vspace{1.6em}
\noindent{\small ［译注：以上前言据底本第 8–9 页译出，只作语言转换，内容未作增删；作者原话的转述保持原意。］}
